\documentclass[11pt]{article}
\usepackage[margin=0.92in]{geometry}
\usepackage{amsmath,amssymb,amsthm,booktabs,array,enumitem,microtype}
\usepackage[T1]{fontenc}
\usepackage{lmodern}
\setlength{\parindent}{0pt}
\setlength{\parskip}{0.48em}
\setlength{\emergencystretch}{3em}
\newtheorem{theorem}{Theorem}
\newtheorem{lemma}{Lemma}
\newtheorem{proposition}{Proposition}
\newcommand{\Sub}{\mathcal{A}}
\newcommand{\PhiSrc}{\Phi_{\mathrm{src}}}
\newcommand{\ModCur}{\mathbf{Mod}_{\mathrm{cur}}}
\newcommand{\IsoCur}{\mathbf{Iso}_{\mathrm{cur}}}
\newcommand{\DiffFour}{\mathbf{Diff}_{4}}
\newcommand{\NoDelta}{\texttt{no\_delta}}
\title{P4--T02 Current-Source Equivariant-Asymmetry\\Derivation Audit}
\author{RT-20260814-002 parent synthesis}
\date{14 August 2026}

\begin{document}
\maketitle

\section{Decisive result}

The exclusive result is
\texttt{current\_source\_equivariant\_asymmetry\_derivation\_theorem},
with identifier
\begin{center}
\small\ttfamily
THM-V22-P4T02-B2-CURRENT-SOURCE-KSTAR-ARROW-SUBCLASS-001.
\end{center}

The maximal exact typed source signature has one semantic primitive:
\(\Sub\) is a smooth four-dimensional manifold carried as primitive debt.
The registered \(\PhiSrc\) token remains an unresolved debt marker rather than
a typed function or relation.  Preserving it that way adds no semantics and
leaves the model groupoid \(\ModCur=\DiffFour\), smooth four-manifolds and all
diffeomorphisms.  Interpreting \(\PhiSrc\) as an action would instead be a
forbidden source extension.

Inside \(\DiffFour\), define
\[
  K_{\!*}(M)=
  \begin{cases}
    \operatorname{Diff}_{c}(M), &
      \operatorname{Diff}_{c}(M)\subsetneq\operatorname{Diff}(M),\\
    \{\operatorname{id}_{M}\}, &
      \operatorname{Diff}_{c}(M)=\operatorname{Diff}(M).
  \end{cases}
\]
Under the ordinary nonempty-manifold convention, this defines an objectwise
proper wide normal arrow subclass, natural under every diffeomorphism and
containing nonidentity arrows globally.  It is determined by the current
smooth topology and factors through the current reduct.  This satisfies the
strongest literal mathematical
\texttt{ASYM-ARROW-SUBCLASS} candidate.  It does not adopt the subclass as
\(\mathrm{Eq}_{\mathrm{src}}\), declare its arrows physically allowed, or
supply a point, root, orientation, order, admission law, causal structure, or
P4--T02 cone selector.

\section{Exact registered signature}

The canonical foundations and dynamics sources give the following exact
formation table.

\begin{center}
\small
\begin{tabular}{>{\raggedright\arraybackslash}p{0.16\linewidth}
                >{\raggedright\arraybackslash}p{0.27\linewidth}
                >{\raggedright\arraybackslash}p{0.43\linewidth}}
\toprule
Symbol & Registered status & Semantic consequence \\
\midrule
\(\Sub\) & Smooth manifold, \(\dim\Sub=4\), primitive debt & A
smooth-manifold reduct has a standard model groupoid. \\
\(\PhiSrc\) & Unresolved abstract source-order or evolution debt marker & No
semantic interpretation is assigned; the typed signature remains the
smooth-manifold primitive. \\
Target benchmark & Exact GR action, metric, coupling, and equations adopted
separately at effective level & None may be used to complete the source
signature. \\
\bottomrule
\end{tabular}
\end{center}

\begin{proposition}[Unresolved-marker boundary]
Representing \(\PhiSrc\) as an uninterpreted debt marker adds no object or
arrow condition to \(\DiffFour\).  Any assignment of a domain, codomain,
arity, action, relation, or preservation equation to it would be a proper
extension of the registered semantic signature.
\end{proposition}

\begin{proof}
The marker supplies no well-formed term to interpret and hence no commutation
square to impose.  The only semantic object type stated by the registered
source package is the smooth four-manifold.  Adding any of the listed data
changes that set of semantic types and equations.
\end{proof}

This reading is exact but modest: it forms only the maximal currently typed
groupoid and does not pretend to resolve the debt marker.

\section{Well-typed bare-reduct controls}

Write \(\DiffFour\) for the groupoid whose objects are smooth
four-manifolds and whose arrows are diffeomorphisms.  Identities, inverses,
and composition preserve exactly the smooth structure and dimension.  The
following calculations concern \(\DiffFour\) only.

\begin{lemma}[Equivariant terminal sections]
If a group \(G\) acts on a set \(P\), then
\[
  \operatorname{Hom}_{G}(1,P)\cong P^{G}.
\]
Thus a natural choice on an object must be fixed by every automorphism of that
object.
\end{lemma}

\begin{proof}
An equivariant map from the singleton is determined by its value \(p\).  Its
equivariance equations are exactly \(g p=p\) for every \(g\in G\).
\end{proof}

\subsection{Point and nonempty finite-set control}

On \(\mathbb{R}^{4}\), let \(T(x)=x+e_{1}\).  Every orbit
\(\{T^{n}p:n\in\mathbb{Z}\}\) is infinite.  Hence no point is fixed by the
translation subgroup.  If a nonempty finite set \(F\) were invariant, it
would contain the infinite orbit of any \(p\in F\), a contradiction.  The
point/finite-set candidate therefore has no natural terminal section on the
bare reduct.

\subsection{Orientation control}

The reflection
\[
 r(x_{1},x_{2},x_{3},x_{4})=(-x_{1},x_{2},x_{3},x_{4})
\]
is a diffeomorphism of \(\mathbb{R}^{4}\) with determinant \(-1\).  It
exchanges the two sheets of the orientation torsor.  Therefore no sheet is
fixed by every automorphism.  This already blocks a natural orientation even
after restricting attention to the orientable control object.  It says
nothing about physical time orientation.

\subsection{Total-order control}

The affine diffeomorphism \(s(x)=e_{1}-x\) exchanges \(0\) and \(e_{1}\).
If a strict total order were invariant under every diffeomorphism, either
\(0<e_{1}\) or \(e_{1}<0\).  Applying \(s\) reverses the selected relation,
contradicting asymmetry of the strict order.  Thus the bare smooth structure
selects no diffeomorphism-invariant total order.

These are full automorphism-subgroup controls on one registered-type model,
not the RT002 two-object pair-groupoid calculation.

\section{Presentation candidates and the positive arrow witness}

Two presentation candidates do not have source-output types in the current
typed signature.
\begin{itemize}[leftmargin=2em]
  \item A root or label needs a presentation carrier, root sort, label sort,
        and transport law.
  \item A nonconstant admission predicate needs a presentation-candidate
        carrier, a Boolean map, and arrow transport.
\end{itemize}

The arrow-subclass candidate is different because compact support is already a
topological notion on every smooth manifold.

\begin{theorem}[Objectwise-proper compact-support arrow subgroupoid]
For every nonempty smooth four-manifold \(M\), define \(K_{\!*}(M)\) by the
displayed two-branch rule.  Let \(\mathcal{C}_{\!*}\) have all objects of
\(\DiffFour\), arrows
\(\mathcal{C}_{\!*}(M,M)=K_{\!*}(M)\), and no arrows between distinct objects.
Then \(\mathcal{C}_{\!*}\) is a wide normal subgroupoid, the assignment is
natural under every diffeomorphism, it is proper at every object, and it
contains nonidentity arrows globally.
\end{theorem}

\begin{proof}
The identity has empty compact support.  If \(f\) has compact support, then
\(\operatorname{supp}(f^{-1})=f(\operatorname{supp}f)\), which is compact.
Moreover
\[
 \operatorname{supp}(f\circ g)
 \subseteq \operatorname{supp}(g)\cup g^{-1}(\operatorname{supp}f),
\]
so composition preserves compact support.  For every diffeomorphism
\(h:M\to N\),
\[
 \operatorname{supp}(hfh^{-1})=h(\operatorname{supp}f),
 \qquad
 h\operatorname{Diff}_{c}(M)h^{-1}=\operatorname{Diff}_{c}(N).
\]
Thus conjugation preserves both \(\operatorname{Diff}_{c}\) and the equality
predicate \(\operatorname{Diff}_{c}=\operatorname{Diff}\), so it transports
\(K_{\!*}(M)\) exactly to \(K_{\!*}(N)\).  Properness holds by definition in
the first branch.  In the second branch, a coordinate ball supports a nonzero
compactly supported smooth vector field whose sufficiently small-time flow is
a nonidentity diffeomorphism; hence \(\operatorname{Diff}(M)\) is nontrivial
and \(\{\operatorname{id}_{M}\}\) is proper.  On \(\mathbb{R}^{4}\), compactly
supported bump flows show the family is not identity-only, while the
translation \(x\mapsto x+e_{1}\) shows
\(\operatorname{Diff}_{c}(\mathbb{R}^{4})\) is proper.
\end{proof}

No cross-object arrow is required by the definition of a wide subgroupoid.
The theorem supplies a mathematical natural arrow subclass, not a normative
law saying that precisely these arrows are source-equivalent or physically
allowed.  That semantic and bridge-facing question remains open.

\section{Current-reduct factorization}

Let \(U:E\to R\) forget proposal-only expansion data and let \(q:E\to Q\)
record one proposed asymmetry.

\begin{theorem}[Fibre-constancy criterion]
There is a map \(\bar q:U(E)\to Q\) such that
\(q=\bar q\circ U\) if and only if \(q\) is constant on every \(U\)-fibre.
The factor is unique on \(U(E)\); uniqueness on all of \(R\) additionally
requires \(U\) to be surjective or values off the image to be specified.
\end{theorem}

\begin{proof}
If \(q=\bar q\circ U\), equal \(U\)-values give equal \(q\)-values.
Conversely, fibre constancy makes \(\bar q(U(e))=q(e)\) well defined.  The
final uniqueness statement is immediate from the domain \(U(E)\).
\end{proof}

For each of the six predeclared candidates, two arbitrary proposal-only
expansions over the same \(\mathbb{R}^{4}\) reduct may carry different
candidate data; those supplied values fail factorization.  The compact-support
witness is different: it is computed from the smooth topology, so its value is
constant on every expansion fibre and does factor.  The six-pair hexad is a
same-current-reduct diagnostic; it neither stipulates a successful root nor
replays the finite rooted-admission fixture.  The executable exact model also
checks all \(2^{6}=64\) Boolean maps on three two-element fibres and finds
exactly \(2^{3}=8\) factorable maps.

\section{Six-candidate audit matrix}

\begin{center}
\small
\begin{tabular}{>{\raggedright\arraybackslash}p{0.19\linewidth}
                >{\raggedright\arraybackslash}p{0.28\linewidth}
                >{\raggedright\arraybackslash}p{0.39\linewidth}}
\toprule
Candidate & Full-package status & Bare-reduct diagnostic \\
\midrule
Point / finite set & Negative naturality & Translation subgroup has no fixed
point or nonempty finite invariant subset. \\
Orientation & Negative naturality & Reflection exchanges orientation sheets. \\
Order & Negative naturality & Affine point exchange contradicts invariant
strict total order. \\
Root / label & Formation failure & No presentation or label type. \\
Admission & Formation failure & No candidate carrier or admission predicate
type. \\
Arrow subclass & Positive theorem & \(M\mapsto K_{*}(M)\)
is natural, objectwise proper, and globally nonidentity. \\
\bottomrule
\end{tabular}
\end{center}

Exactly one candidate is positive.  The result follows from smooth topology
alone and leaves \(\PhiSrc\) unresolved.  The five negative candidates remain
useful boundaries on what the theorem does not select.

\section{Result classification and continuation boundary}

The current-source derivation branch is selected because the compact-support
assignment is typed, natural under every current arrow, and current-reduct
factorable.  Negative fixed-point and formation branches remain candidate-local
and do not replace the exclusive positive result.  No target or desired-output
premise was used, and no source-extension candidate was constructed.

The separately admitted future packet is
\begin{center}
\small\ttfamily
PKT-V22-P4T02-B2-POST-CURRENT-SOURCE-EQUIVARIANT-ASYMMETRY-\\
DERIVATION-THEOREM-THEORETICAL-CONTINUATION-SELECTION-V1.
\end{center}
Its type is \texttt{theoretical\_continuation\_selector}, its role is
Theoretical Continuation Selector version 0.1.0, and it is selected but not
executed.  It must decide whether the compact-support theorem has any
bridge-facing P4--T02 relevance or is source-side classification only.  Route C
is not preselected here.

\section{Freeze and Distance-to-GR status}

All six inherited local freezes remain unchanged:
\begin{enumerate}[leftmargin=2em]
  \item \texttt{NDCL-V22-P4T02-B2-SHARED-TAU-SELECTOR};
  \item \texttt{NDCL-V22-P4T02-B2-REPAIRED-QUOTIENT-DESCENT-ROBUSTNESS};
  \item \texttt{NDCL-V22-P4T02-B2-COMMON-CHARACTER-HOLONOMY-PROTECTION};
  \item \texttt{NDCL-V22-P4T02-B2-ORIENTED-MATROID-BRIDGE-SELECTION-ROBUSTNESS};
  \item \texttt{NDCL-V22-P4T02-B2-SIGNED-CUBIC-VIABILITY-SELECTOR-ROBUSTNESS};
  \item \texttt{NDCL-V22-P4T02-B2-SOURCE-LAW-SPACE-ROBUST-INVARIANCE-}
        \texttt{PROTECTION-ROBUSTNESS}.
\end{enumerate}

Every expanded Distance-to-GR burden remains \NoDelta{}: source ontology
primitives, source equivalence, RetainH, GenH, ObsLoc, Resp, source manifold,
effective metric, matter coupling, Einstein equations, finite-variation
robustness, benchmark promotion, Gate Chair status, and current route-freeze
status.  The result adds a mathematical classifier but no physical object.

\section{Non-conclusions}

The current model groupoid and compact-support subgroupoid are not adopted
\(\mathrm{Eq}_{\mathrm{src}}\), physical gauge groups, allowed physical-arrow
laws, or causal structures.  Mathematical orientation is not physical time
orientation.  No source primitive, source law, point/root selector,
presentation ontology, empirical response, universal propagation, effective
metric, matter coupling, or Einstein equation is derived.  D7 is not
reevaluated, B2 remains inactive, and P4--T03 remains locked.  The theorem does
not imply that every conservative source extension is unnecessary or
impossible.  No Gate verdict, benchmark promotion, proof authority,
publication, push, external action, or completed derivation follows.

\end{document}
